\section{Conditional half-counts and the certificate iteration}
\label{p14:comp:half}

We prove Theorem~\ref{p13:half:conditional} with its stated uniformity
in arbitrary fixed source blocks.  Blocks and effective spaces have the
meaning specified in Section~\ref{p13:half:section}; in particular,
the nonzero projections are required in the individual new spaces,
not in their quotients by the old space.

\subsection{Small kernels and smooth moduli}

Choose \(H=C_DB\) larger than every block diameter, and set
\[
 K_H(n)=\prod_{\substack{\ell>H\\v_\ell(n)\ {\rm odd}}}\ell.
\]
For \(q>0\), let \(\mathcal X_q\) be the set of centres
containing distinct values \(n_1,n_2\) with
\(K_H(n_1)K_H(n_2)\le N^q\).

\begin{lemma}[Two small kernels]
\label{p13:half:small-kernels}
For every fixed \(q>0\) and \(\varepsilon>0\),
\(\#\mathcal X_q\ll_{q,D,C,\varepsilon}N^{q+\varepsilon}\).
\end{lemma}

\begin{proof}
Write \(n_i=c_i k_i a_i^2\), where \(k_i=K_H(n_i)\)
and \(c_i\) is squarefree and supported on primes at most \(H\).
The \(c_i\) have \(2^{2\pi(H)}=N^{o(1)}\) choices, while
the pairs \(k_1k_2\le N^q\) have \(N^{q+o(1)}\) choices.
Fix the two distinct offsets. Their nonzero difference \(h\)
gives
\[
 c_1k_1a_1^2-c_2k_2a_2^2=h.
\]
The effective coefficients are \(O_D(N)\).
Lemma~\ref{p1:arith:pell} gives \(N^{o(1)}\) solutions,
including the factorizable case, and each solution determines
the centre. The \(O_D(B^2)\) offsets are absorbed.
\end{proof}

For \(Y=(\log N)^A\), with fixed \(A\ge2\), write
\[
 n=s_Y(n)a_Y(n)^2K_Y(n).
\]
Here \(s_Y(n)\) is the entire part supported on primes at most
\(Y\); \(a_Y(n)\) contains the halves of the valuations at
larger primes; and \(K_Y(n)\) is their odd squarefree part.

\begin{lemma}[A large smooth modulus]
\label{p13:half:smooth}
Fix \(t\), \(c>0\), and \(A\ge2\).
There is a deterministic set \(\mathcal M_Y\) of centres,
of cardinality \(O(N^{1-1/A+\varepsilon})\), with the
following property. If a selection has more than
\(c\log\log N\) values, its centre lies outside
\(\mathcal X_{1/2}\), and it contains fewer than \(t\)
values with \(K_Y(n)>N^{1/5}\), then its centre belongs
to \(\mathcal M_Y\).
\end{lemma}

\begin{proof}
Outside \(\mathcal X_{1/2}\), at most one value satisfies
\(a_Y(n)^2>N^{3/4+1/100}\): two such values would give
\[
 K_H(n_1)K_H(n_2)
 \le\frac{O_D(N^2)}{a_Y(n_1)^2a_Y(n_2)^2}<N^{1/2}.
\]
Remove this possible value and the at most \(t-1\)
large kernels. For sufficiently large \(N\), select
26 remaining values. Their smooth parts satisfy
\[
 s_i\gg_D N^{1-3/4-1/100-1/5}=N^{1/25}.
\]
Since \(\gcd(s_i,s_j)\le |n_i-n_j|=O_D(B)\),
\[
 \operatorname{lcm}(s_1,\ldots,s_{26})
 \ge\frac{\prod_i s_i}{\prod_{i<j}\gcd(s_i,s_j)}
 \ge N^{26/25-o(1)}>C_1N,
\]
where \(C_1N\) exceeds the length of the centre interval.
Multiply its prime factors, one at a time and with
multiplicity, until the first passage above \(C_1N\).
The resulting smooth modulus satisfies
\[
 C_1N<D_0\le C_1NY,\qquad
 D_0\mid\operatorname{lcm}(s_1,\ldots,s_{26}).
\]
For fixed \(1-1/A<\sigma<1\), Rankin's elementary bound
counts such moduli by
\[
 (C_1NY)^\sigma
       \prod_{\ell\le Y}(1-\ell^{-\sigma})^{-1}
 =N^{\sigma+o(1)}.
\]
Indeed the logarithm of the product is
\(O_\sigma(Y^{1-\sigma})=o(\log N)\).
Assign each maximal prime power of \(D_0\) to one of the
26 offsets on which it occurs. This costs
\(26^{\omega(D_0)}=N^{o(1)}\).
The Chinese remainder theorem determines at most one centre
in its interval, since \(D_0>C_1N\). The offsets themselves
have \(B^{26}=N^{o(1)}\) choices.
Choose \(\sigma\) sufficiently close to \(1-1/A\).
\end{proof}

\subsection{Several anchors force a congruence certificate}

\begin{lemma}[Incoming and outgoing certificate counts]
\label{p13:half:saturation}
Suppose a relation in \eqref{p13:half:condition} selects
\(t\) values in one new block whose kernels exceed \(N^b\),
where \(b=1/5\) and \(tb>r+s-1\).
Some other old or new block receives a congruence certificate
whose modulus exceeds a fixed multiple of \(NB\).
For fixed \(r,s,t,A\), the directed graph of these
source--receiver certificates has both incoming and outgoing
degrees at most \(N^{1/A+o(1)}\).
\end{lemma}

\begin{proof}
Represent the old vector by a selection of actual old values;
parity in the old blocks is not required. Together with the
new selections this is a square product.
Primes above \(Y\) cannot occur in two distinct values
of one block. The \(t\) anchor kernels are therefore
pairwise coprime, and their product \(Q\) satisfies
\(Q>N^{tb}\), \(Q\ll_D N^t\).
Assign each prime of \(Q\) to one of its other odd occurrences.
The \(r+s-1\) receiver moduli have product \(Q\).
For sufficiently large \(N\), at least one exceeds \(C_2NB\).

In that modulus keep the primes only up to its first passage
above \(C_2NB\). The retained list has at most
\[
 k=\left\lfloor\frac{\log(C_2NB)}{\log Y}\right\rfloor+1
\]
members. For a fixed source, the \(t\) offsets and a subset
of the \(O_t(\log N/\log Y)\) available primes cost
\(N^{o(1)}\). Assigning their destinations among \(O_D(B)\)
receiver positions costs at most
\(B^k=N^{1/A+o(1)}\). The congruences determine the receiver.

For a fixed receiver, its values contain
\(R=O_D(B\log N/\log Y)\) primes above \(Y\).
Choose at most \(k\) of them and assign each to an anchor:
the number of codes is at most
\[
 B^t\sum_{j\le k}\binom Rj t^j=N^{1/A+o(1)}.
\]
When \(R\ge2k\), use \((k+1)(etR/k)^k\);
when \(R<2k\), use \((1+t)^R=N^{o(1)}\).
The retained modulus again exceeds the source interval.
Forgetting all other primes enlarges the count and does not
remove the original full-support relation from the remaining
blocks.
\end{proof}

\subsection{The iteration}
\label{p14:comp:half:iteration}

\begin{proof}[Proof of Theorem~\ref{p13:half:conditional}]
We induct strongly on \(r\), simultaneously for every fixed
number \(s\) of sources. Put \(T_0=1\). The case \(r=1,s=0\)
is empty, since its sole projection is required to be nonzero.
For all other cases, choose
\[
 t=5(r+s-1)+1.
\]
First choose a small fixed step \(\delta>0\) within the
final error margin. Let
\[
 x_0=0,\qquad x_j=\min\{\tfrac12,x_{j-1}+\delta\},
 \qquad A_j=x_j^{-1},\qquad Y_j=(\log N)^{A_j}.
\]
There are \(J=\lceil1/(2\delta)\rceil\) stages.
All these numbers are fixed before \(N\); choose the
internal margins after \(J\).

Corollary~\ref{p13:short:old-space} counts lists possessing
a witness with all new selections short.
For every remaining list choose one witness, without
enumerating it. It has at least one long new selection.
At the start of stage \(j\), every surviving long new
source outside \(\mathcal X_{1/2}\) belongs to a set
of size \(N^{1-x_{j-1}+\varepsilon'}\).
This is trivial for \(j=1\).

If such a source selects \(t\) large \(K_{Y_j}\) kernels,
Lemma~\ref{p13:half:saturation} supplies a receiver.
When the receiver is new, count the source--receiver pair
by \(N^{1-x_{j-1}+x_j+O(\varepsilon')}\).
The other \(r-2\) projections are still nonzero in their
own spaces; their sum belongs to the old space enlarged
by these two fixed blocks. The induction hypothesis counts
them by \(N^{(r-2)/2+\varepsilon'}\).
The total is
\[
 N^{r/2+x_j-x_{j-1}+O(\varepsilon')}
 \le N^{r/2+\delta+O(\varepsilon')}.
\]
If the receiver is old, it is already fixed. The incoming
degree counts the new source by \(N^{x_j+O(\varepsilon')}\).
Adjoin that source to the old space and apply induction
to the \(r-1\) remaining targets. The cost is at most
\[
 N^{x_j+(r-1)/2+O(\varepsilon')}
 \le N^{r/2+O(\varepsilon')}.
\]
Only a fixed number of choices of roles is involved.

Remove all these saturated cases. Every surviving long
selection outside \(\mathcal X_{1/2}\) has fewer than \(t\)
large kernels, and Lemma~\ref{p13:half:smooth} places its
centre in a set of size \(N^{1-x_j+\varepsilon'}\).
This proves the stage invariant.
At the final stage every long source lies in
\(\mathcal X_{1/2}\cup\mathcal M_{Y_J}\), whose size is
\(N^{1/2+\varepsilon'}\).
Fix one such centre and use the induction hypothesis for
the other \(r-1\) targets. This costs
\(N^{1/2+\varepsilon'}N^{(r-1)/2+\varepsilon'}\).

Every application of induction uses fewer targets,
although the fixed number of old sources can increase.
Thus the simultaneous induction is legitimate. Choosing
\(\delta\) and then all internal margins within
\(\varepsilon\) proves the theorem.
\end{proof}
